
\subsection{Ortogonalių matricų grupė \texorpdfstring{$\Orth{3}$}{O(3)}}


Tegul $(\cdot,\cdot) : \mathbb{R}^3 \times \mathbb{R}^3 \to \mathbb{R}$
būna standartinė skaliarinė sandauga.

\begin{definition}
	Ortogonalių matricų grupė $\Orth3$ yra sudaryta iš $3\times3$ matricų
	$R$ virš realiųjų skaičių, tokių kad, lygybė 
	\begin{equation}
		(R \vb v, R \vb w) = (\vb v, \vb w)
	\end{equation}
	galioja visiems vektoriams $\vb v, \vb w \in \mathbb{R}^3$. 
\end{definition}

Matricos $R \in \Orth3$ vadinamos ortogonaliomis,
nes bet kokie ortogonalūs vektoriai $\vb v$ ir $\vb w$
yra atvaizduojami į ortogonalius vektorius $R\vb v$ ir $R\vb w$.
Šios matricos taip pat pasižymi tuo,
kad matricą $R$ sudauginus su savo transponuotą $R^t$
rezultatas yra vienetinė matrica $I$ 
\begin{equation}
	\begin{split}
		\det(RR^t) = \det(I)
		\\
		\qty(\det(R))^2 = 1
		\\
		\det(R) = \pm 1
	\end{split}
\end{equation}
Determinantas yra homomorfizmas iš $\Orth3$ į grupę $\qty{1,-1} \cong C_2$.
To pasekoje $\Orth3$ Yra padalijama į dvi dalis:
tikruosius sukimus (\textit{ang.} proper rotations),
kuriems $\det(R) = 1$,
ir netikruosius sukimus (\textit{ang.} improper rotations),
kuriems $\det(R) = -1$.
\begin{example}
	Sukimas kampu $\alpha$ aplink $z$ ašį užrašomas matrica
	\begin{equation*}
		\mqty(
		\cos(\alpha) & -\sin(\alpha) &   \\
		\sin(\alpha) &  \cos(\alpha) &   \\
			     &		     & 1
		)
	\end{equation*}
	netikrasis sukimas,
	kuris vektorių atspindi per $xy$ plokštumą,
	užrašomas matrica
	\begin{equation*}
		\mqty(
		1 &   &   \\
		  & 1 &   \\
		  &   & -1
		) 
	\end{equation*}
	atspindys per tašką $0$ yra 
	netikrasis sukimas vadinamas inversija,
	jis užrašomas matrica
	\begin{equation*}
		P :=\mqty(
		-1 &    &   \\
		   & -1 &   \\
		   &    & -1
		) 
	\end{equation*}
\end{example}


Tikrieji sukimai yra fiziškai realizuojami,
kaip objekto krypties pakeitimas erdvėje.
Netikrieji sukimai nėra fiziškai realizuojami,
jiems atlikti reikia objektą pakeisti savo atspindžiu
ir tada pritaikyti tikrąjį pasukimą.
Tikrieji sukimai sudaro grupės $\Orth3$ pogrupį vadinama $\SO3$.

\begin{definition}
	Specialioji ortogonalių matricų grupė žymima $\SO3$ ir yra sudaryta iš
	$3\times3$ matricų $R$ virš realiųjų skaičių, tokių kad,
	\begin{itemize}
		\item lygybė $(R \vb v, R \vb w) = (\vb v, \vb w)$
			galioja visiems vektoriams 
			$\vb v$, $\vb w \in \set R^3$ 
		\item ir $\det(R) = 1$.
	\end{itemize}
\end{definition}


\begin{proposition}
	Grupė $\Orth3$ yra izomorfiška grupių sandaugai
	$\SO3 \times \qty{1,-1}$.
\end{proposition}
\begin{proof}
	Tegul $f$ yra atvaizdavimas
	\begin{equation}
		\begin{aligned}
			f: \SO3 \times \qty{1,-1} &\to \Orth3\,,\\
			(R,1) &\mapsto R \\
			(R,-1) &\mapsto PR \,.
		\end{aligned}
	\end{equation}
	Įrodome, kad $f$ atvirkštinis atvaizdavimas yra
	\begin{equation}
		\begin{alignedat}{4}
			f^{-1}: \Orth3 &\to \SO3 \times \qty{1,-1}
			\,,\\
			R &\mapsto (R,1)  &R&\in \SO3 \\
			R &\mapsto (PR,-1)  &R&\notin\SO3
			\,.
		\end{alignedat}
	\end{equation}
	Turime $f \circ f^{-1} = \mathup{id}$
	\begin{equation}
		\begin{aligned}
			\text{kai } R \in\SO3 
			\,,\quad
			&f\circ f^{-1} (R) = f((R,1)) = R
			\,,\\
			\text{kai } R \notin\SO3 
			\,,\quad
			&f\circ f^{-1} (R) = f((PR,-1)) = PPR = R
			\,.
		\end{aligned}
	\end{equation}
	Turime $f^{-1} \circ f = \mathup{id}$
	\begin{equation}
		\begin{aligned}
			&f^{-1}\circ f ((R,1)) = f^{-1}(R) = (R,1)
			\,,\\
			&f^{-1}\circ f ((R,-1)) = f^{-1}(PR) = (PPR,-1)= (R,-1)
			\,.
		\end{aligned}
	\end{equation}
	Kad įrodyti kad $f$ yra grupių izomorfizmas lieka įrodyti,
	kad $f$ išlaiko grupės struktūrą dauginant teisingus sukimus,
	teisingą sukimą su neteisingu ir dauginant du neteisingus sukimus.
	\begin{equation}
		\begin{split}
			f((R_1,1))f((R_2,1))
			&= R_1 R_2 \\
			&= f((R_1 R_2,1)) \\
			&= f((R_1,1)(R_2,1)) \,,
		\end{split}
	\end{equation}
	\begin{equation}
		\begin{split}
			f((R_1,1))f((R_2,-1))
			&= R_1 P R_2 \\
			&= P R_1 R_2 \\
			&= f((R_1 R_2,-1)) \\
			&= f((R_1,1)(R_2,-1)) \,,
		\end{split}
	\end{equation}
	\begin{equation}
		\begin{split}
			f((R_1,-1))f((R_2,-1))
			&= P R_1 P R_2 \\
			&= P^2 R_1 R_2 \\
			&= R_1 R_2 \\
			&= f((R_1 R_2,1)) \\
			&= f((R_1,-1)(R_2,-1)) \,.
		\end{split}
	\end{equation}
	Čia pasinaudojome faktu kad $PR=RP$ su visais 
	$R \in \SO3$ ir $P^2=I$. 
	
	Egzistuoja grupių izomorfizmas $f$,
	tad $\Orth3$ ir $\SO3 \times \qty{1,-1}$
	yra izomorfiškos grupės.
\end{proof}




